\documentclass[a4paper]{resume} % Use the custom resume.cls style
\usepackage{nopageno}
\usepackage[left=0.25in,top=0.15in,right=0.15in,bottom=0.25in]{geometry} % Document margins
\usepackage{fontawesome}
\usepackage{times}
    \newcommand{\tab}[1]{\hspace{.2667\textwidth}\rlap{#1}}
    \newcommand{\itab}[1]{\hspace{0em}\rlap{#1}}
\usepackage[german]{babel}
\usepackage{ulem}
\usepackage{enumitem}
    \setitemize{noitemsep,topsep=0pt,parsep=0pt,partopsep=0pt}
\usepackage{graphicx}

\name{José Alfredo Rosas Córdova} % Your name 
\address{\faEnvelope{ jose.alfredo.rosas.cordova@stud.uni-hannover.de}\hspace{0.75em}\faPhone{ +49 162 5197979}}
\address{\faMapMarker{ Heidjerhof 3, Zi. 103, 30625 Hannover}} % Your address
\address{\faBirthdayCake{ 22.10.1999}\hspace{0.75em}\faGlobe{ Guadalajara, Mexiko}}

\begin{document}
    % === AUSBILDUNG ===
    \begin{rSection}{Ausbildung}

        {\bf Maschinenbau, M. Sc.} \hfill {\itshape 10.2025 --- heute} 
        \\{ \textit{Leibniz Universität Hannover}}
        %\\ Notendurchschnitt: 97.97 (entspricht 1.2)
        
        {\bf Luftfahrttechnik, B. Eng.} \hfill {\itshape 08.2017 --- 06.2023} 
        \\{ \textit{Universidad Autónoma de Nuevo León}}
        \\ Abschluss mit summa cum laude
        %\\ Notendurchschnitt: 97.97 (entspricht 1.2)

        {\bf Kombinierte Studien- und Praxisaufenthalte (DAAD-KOSPIE)} \hfill {\itshape 09.2021 --- 09.2022} 
        \\{\bf Maschinenbau, B. Sc.}
        \\ { \textit {Technische Universität Braunschweig}}
        \\ Notendurchschnitt: 1.9
    \end{rSection}
    
    \begin{rSection}{Sprachen}
        \begin{tabular}{ @{} >{\bfseries}l @{\hspace{6ex}} l }
            Spanisch & Muttersprachler\\
            Englisch & TOEFL ITP -- 677 Punkte (C1)\\
            Deutsch & TestDaF -- TDN 5 (C1)
        \end{tabular}
    \end{rSection}
    
    \begin{rSection}{IT-Kenntnisse}
        \begin{tabular}{ @{} >{\bfseries}l @{\hspace{6ex}} l }
            CAD \& CAM: & PTC Creo, OnShape, SOLIDWORKS, Fusion 360 \\
            Programmierkenntnisse: & Python, MATLAB\\
	        Simulation \& Modellierung: & ANSYS Mechanical, Fluent, Simulink, Xfoil\\
            Andere: & MS Office, Oracle PLM \& ERP, \LaTeX, Git\\
        \end{tabular}
    \end{rSection}

    % ==== WORK EXPERIENCE SECTION ====
	\begin{rSection}{Berufserfahrung}
	
		{\bf Wissenschaftliche Hilfskraft} \hfill{\itshape 04.2026 --- heute} 
		\\{\textit{\textbf{Institut für Turbomaschinen und Fluid-Dynamik} -- Hannover, Deutschland}}
		\begin{itemize}
			\item Nutzung spezialisierter Software zur Extraktion von Verformungsinformationen einer Windenergieanlage mit markierten Schaufeln mittels Digital Image Correlation-Techniken. 
			\item Nachbearbeitung der erhaltenen großvolumigen Datensätze mit selbst- und weiterentwickelten Skripten mit Python und MATLAB, die stark auf parallele Programmierschemata angewiesen sind. Die Verarbeitungsschritte umfassten unter anderem die Klassifikation, Transformation und Abbildung der gemessenen Punkte sowie die Erstellung der für die folgenden Validierungsschritte benötigten Diagramme. 
			\item: Intensiver Einsatz von Git-basierten Tools zur Verwaltung von Softwareversionen und zur sicheren Implementierung neuer Funktionalitäten.
			\item Verwendung eines Blender-basierten digitalen Zwillings, um die Richtigkeit der implementierten Algorithmen zu validieren, bevor sie zur Auswertung realer experimenteller Daten verwendet werden.
		\end{itemize}
	
	
        {\bf Mechanical Engineer} \hfill{\itshape 07.2023 --- 09.2025} 
        \\{\textit{Entwicklung \& Einführung neuer Produkte, Hydronische Kühlung}}
        \\{\textit{\textbf{Vertiv} -- Monterrey, Mexico}}
        \begin{itemize}
            \item Entwicklung, Dokumentation und Produktpflege mechanischer Komponenten und Baugruppen wie z.B. Blechschweißkonstruktionen; Primär- und Sekundärkühlkreisläufe (d. h. Rohrleitungsnetze für Kältemittel und Wasser, aus Kupfer bzw. Edelstahl) für neue Produkte, die für hohe Dichte Kühlanwendungen in Rechenzentren gedacht sind (z. B. Kältemaschinen, CDUs, Fluidnetzwerke). Alles in einer kollaborativen Arbeitsumgebung, die auf dem Einsatz von Oracle PLM- und ERP-Systemen basiert ist.
            \item Zusammenarbeit mit Fertigungs- und Außendienstteams und -einrichtungen für die Implementierung von Designverbesserungen hinsichtlich DFM\&A und Wartungsfreundlichkeit sowohl in Design- als auch in Produktionsumgebungen.
            \item Teilnahme an der Planung, der Teilebeschaffung und dem Bauprozess sowohl des Engineerings als auch der Manufacturing-Prototypen brandneuer Produkte sowohl in Mexiko als in den USA.
            \item Verarbeitung von Testdaten zur Designvalidierung für die Überprüfung der Leistung neuer Produkte.
        \end{itemize}
        
        {\bf Project Engineer - UAV Development} \hfill{\itshape 01.2023 --- 04.2023} 
        \\{\textit{\textbf{Spacelab MX} -- Monterrey, Mexico}}
        \begin{itemize}
            \item Bewertung der Durchführbarkeit des Missionskonzepts und dessen Verbesserung.
            \item Erste Flugzeugdimensionierung und vorläufige Flugleistungsstudien.
            \item Entwurf der ersten Flugzeugzellenstruktur für die Prototypenphase des Projekts.
            \item Integration, Boden- und Flugerprobung des Prüfstandsflugzeugs.
        \end{itemize}
    
        {\bf System Testing \& Mechanical Design Intern, Engineering Mechatronics Department} \hfill{\itshape 04.2022 --- 08.2022} 
        \\{\textit{\textbf{Bosch Rexroth AG} -- Stuttgart, Deutschland}}
        \begin{itemize}
            \item Konzeption und Durchführung von Systemtests innerhalb einer Produktionsumgebung für autonome mobile Roboter (UGV) auf Basis zuvor definierter KPIs, die rund um die Kundenanforderungen von Beta-Testern etabliert wurden.
            \item Bewertung, Analyse und Fehlererkennung der Antriebsleistung der UGVs sowie der Kommunikationsschicht zwischen UGVs und Flottenmanagement-Service auf Basis von Testdaten.
            \item Handhabung und Wartung von autonomen mobilen Robotern.
            \item Mechanische Konstruktion von internen Hardwareelementen und Antriebsmechanismen.
            \item Prototyping von vorgeschlagenen Lösungen für das Komponentendesign neben der Validierung im Entwicklungskontext.
        \end{itemize}
    \end{rSection}

    
    % === PROJECTS AND SEMINARS ===
    \begin{rSection}{Forschungserfahrung}
        {\bf Studentische Hilfskraft, Aerodynamik-Labor} \hfill {\itshape 07.2020 --- 06.2023} 
        \\{\textit{Centro de Investigación e Innovación en Ingeniería Aeronáutica in Mexiko}}
        
        \begin{itemize}
        	
	         \item Mitveröffentlichung eines Artikels zur Entwicklung, Manufaktur und Flugversuche meines Abschlussprojekts, wobei Faserverbundwerkstoffe ein wesentlicher Teil der Strukturkomponenten des Flugzeugs sind. Der Artikel wurde bei \textit{Aircraft Engineering and Aerospace Technology} eingereicht.\hfill {\itshape 09/2025}.		
              \begin{itemize}
	         	\item \textit{Design, Manufacture, and Testing of a Low-Cost Tilt-Rotor UAV: A Case Study}. (DOI: 10.1108/AEAT-05-2025-0168). 
	         \end{itemize}

            \item Mitveröffentlichung eines Artikels zur Validierung von Potenzialströmungslösern in der Fachzeitschrift \textit{Revista Internacional de Métodos Numéricos para Cálculo y Diseño en Ingeniería}.\hfill {\itshape 11/2024}
            \begin{itemize}
                \item \textit{Validation of VSPAERO for Basic Wing Simulation} (DOI: 10.23967/j.rimni.2024.10.56492). 
            \end{itemize}
        \end{itemize}
    \end{rSection}
    

    
    % === Extracurricular Section ===
    \begin{rSection}{Extracurriculare Aktivitäten}
      	{\bf Verein Mitglied} \hfill {\itshape 10/2025 --- today} 
    	\\{\textit{\textbf{Akademischer Fliegergruppe Hannover e.V.}}}
    	\vspace{-0.5em}
    	\begin{itemize}
    		\item Mitmachen beim Wartund und Instandhaltund der Räumlichkeiten und Flugzeuge des Vereins.
    	\end{itemize}
    	
    	
        {\bf Team Mitglied (u.a. Verfasser von Technischen Berichten, Leitender Ingenieur)} \hfill {\itshape 09.2019 --- 06.2022} 
        \\{\textit{Axios Aero Design}}
        
        \begin{itemize}
        	\item Erstellung eines technischen Berichts, der aus einer Dokumentation des Konstruktionsprozesses sowie aus technischen Zeichnungen besteht, die das Endprodukt beschreiben, das innerhalb SAE Design-Bettweberbe presentiert wurden.
            \item Konzept- und Vorentwurfsstudien von Flugzeugen, die Aerodynamik, Strukturen, elektrische Antriebe sowie Flugstabilität und -steuerung umfassen. 
            \item Projektplanung und -ausführung auf der Grundlage von System-Engineering-Methoden und Unterstützung bei der Detailkonstruktion von Flugzeugstrukturkomponenten wie Flügel-, Rumpf- und Leitwerksstrukturen.
            \item Erste Praxiserfahrung mit der Herstellung von 3D-Gedrückten Teilen und Komponenten aus Verbundwerkstoffen. Validierung ihrer Umsetzung als Strukturtelekomponenten in einem Wettbewerb, in dem Balsaholzflugzeuge die Regel waren.
        \end{itemize}
    \end{rSection}
    
    % === PRICES ===
    \begin{rSection}{Auszeichnungen}
        {\bf SAE Aero Design México 2022} \hfill {\itshape 03/2022}
        \\{\textit{2. Platz (Gesamt)}}
        \\{\textit{3. Platz (Technischer Bericht)}}
    
        {\bf MathWorks Minidrone Competition Mexico 2021} \hfill {\itshape 08/2021}
        \\{\textit{1. Platz}}
    \end{rSection}

    % === SIGNATURE ===
    \begin{rSection}{}
        \vspace{5em}
        \begin{minipage}{0.5\linewidth}
            \uline{Hannover \today}\\
            \textit{Ort, Datum}
        \end{minipage}
        \begin{minipage}{0.5\linewidth}
            \uline{%
            	José Alfredo Rosas Córdova%
            	\hspace*{0.0em}\rlap{\raisebox{-1.75em}[0pt][0pt]{\includegraphics[height=7em]{/home/jarc/Nextcloud/100 Personal/110 Documents/Recursos Gráficos/blue_signature.png}}}%
            	\hspace{5em}%
            }\\
            \textit{Unterschrift}
        \end{minipage}
    \end{rSection}
\end{document}